% ============================================================
%  CHAPTER V — CONCLUSIONS AND RECOMMENDATIONS
% ============================================================

\chapter{Conclusions and Recommendations}

\section{Conclusions}

This research successfully built a Python-based pipeline that 
integrates automated legacy PHP code conversion, security 
scanning, and ISO/IEC 27001:2022 compliance mapping into a 
single workflow that runs entirely in a local environment. Based 
on testing across two real-world case studies, the following 
conclusions address the research objectives.

\textbf{First}, Rector-based automated PHP conversion was 
successfully implemented, achieving a success rate of 97.6\% on 
the FT UGM dataset and 91.7\% on the CBT DTI UGM dataset, with 
post-conversion syntax validity reaching 100\% in both cases. 
Failures in two files across both datasets were caused by an 
internal Rector bug, not by the source code itself.

\textbf{Second}, the security checking mechanism mapped to 
ISO/IEC 27001:2022 Annex A controls was implemented in a 
deterministic and fully auditable manner. The value $\Delta F = 0$ 
in both case studies confirms that Rector's conversion introduced 
no new vulnerabilities into the codebase --- an important 
confirmation given that large-scale automated transformation has 
the potential to alter code behavior in unexpected ways. The 
resulting \texttt{NON\_COMPLIANT} status is not a reflection of 
pipeline failure, but rather an accurate portrait of the legacy 
code's actual condition, which already contained security issues 
prior to remediation. All analysis results are consolidated into 
a JSON report generated automatically on each pipeline execution, 
making it directly usable as audit documentation without 
additional manual work --- addressing the real need of 
organizations undergoing certification that require written 
evidence that code migration was carried out in accordance with 
applicable security standards.

\textbf{Third}, the PHPStan validation component delivers 
concrete value by exposing thousands of potential type 
compatibility issues that would otherwise go undetected until 
the code is actually run in a production environment. The 
majority of findings represent technical debt that existed long 
before migration and became visible only because PHP 8.3's type 
system is far stricter than its predecessors. This fact 
underscores the pipeline's added value: it does not merely 
convert syntax, but also provides a comprehensive picture of the 
actual state of the codebase that was previously hidden behind 
the leniency of the type system in older PHP versions.

\textbf{Fourth}, the performance comparison of five open-source 
LLM models in the context of PHP code security analysis yields 
empirical findings relevant to anyone seeking to apply a similar 
approach. Format Compliance Rate and Syntax Validity Rate proved 
to be two independent, uncorrelated dimensions: a model that 
faithfully follows format instructions does not necessarily 
produce syntactically valid code, and vice versa. Qwen2.5-Coder 
is recommended as the primary model because it consistently 
achieved an SVR of 100\% across all experimental conditions, 
even while completely ignoring format instructions, whereas other 
models with higher format compliance more frequently produced 
code that failed syntax validation. SVR proved to be the most 
valid evaluation metric for locally running LLM-based systems 
because its results can be independently verified using 
\texttt{php -l}, unlike FCR --- which is susceptible to minor 
response format variations --- or self-reported confidence 
scores, which showed no correlation with the actual quality of 
the generated code.

\section{Recommendations}

Based on the identified limitations, the following development 
directions are recommended for future research.

\begin{enumerate}
    \item LLM evaluation needs to be strengthened from two 
    directions: semantic validation using automated test suites 
    to verify the logical correctness of remediation suggestions, 
    and dataset expansion with a greater diversity of 
    vulnerability types so that model comparison conclusions can 
    be more broadly generalized. As a long-term continuation, 
    fine-tuning exploration with LoRA or QLoRA can be pursued 
    using this research as a baseline.

    \item The accuracy of PHPStan reports on CodeIgniter 3.x 
    codebases can be improved through custom stubs or PHPStan 
    baseline configuration, so that only findings introduced by 
    the conversion are reported.

    \item Validation of the ISO/IEC 27001:2022 mapping by 
    certified auditors or practitioners should be conducted to 
    strengthen the credibility of the generated compliance 
    reports.

    \item Integration into CI/CD pipelines and extension to 
    other programming languages are worth exploring. The 
    pipeline's modular architecture facilitates the replacement 
    of the Rector and Semgrep components with equivalent tools 
    for Python, JavaScript, or other languages, while also 
    opening the possibility of integration with GitHub Actions 
    or GitLab CI.

    \item Testing against PHP 8.4 and 8.5 targets and 
    comparison with commercial tools such as SonarQube or Snyk 
    on the same dataset should be conducted to validate the 
    pipeline's coverage and accuracy more comprehensively.
\end{enumerate}